\subsection*{Plan de gestión del proyecto}
\subsubsection*{Alcance}

La investigación adopta un enfoque tecnológico aplicado con implementación incremental mediante \textbf{Scrum}. Esta decisión responde a la necesidad de construir una solución funcional, validarla con evidencia de uso y ajustar el producto de forma iterativa hasta alcanzar una versión estable orientada a reducir tiempos de espera en transporte público urbano.

El alcance del proyecto comprende: (i) diseño metodológico, (ii) levantamiento de requerimientos, (iii) implementación de la propuesta web con integración GPS, (iv) validación mediante métricas operativas y (v) análisis de resultados para sustentar conclusiones. No se contempla una implementación masiva a escala ciudad; la validación se desarrolla en entorno piloto controlado.

\subsubsection*{Metodología de desarrollo (Scrum)}

Scrum se selecciona por tres razones principales: (i) permite trabajar con entregables cortos y verificables, (ii) facilita priorización por valor en cada ciclo, y (iii) mantiene una retroalimentación continua entre diseño, desarrollo y validación de resultados.

\subsection*{Sujetos/unidades de estudio}

Las unidades de estudio son los usuarios del transporte público urbano de Arequipa y los actores operativos vinculados al servicio (conductores y entidades de gestión). El análisis se centra en la interacción usuario--ruta--tiempo, considerando como variable principal el tiempo de espera y desplazamiento antes y después del uso del sistema web con apoyo de datos GPS.

Como unidades técnicas se consideran los registros de ubicación y eventos operativos (hora de consulta, tiempo estimado de llegada, ruta sugerida, confirmación de viaje y tiempo efectivo de desplazamiento), por constituir la base empírica de la validación.

\subsection*{Universo, población y muestra}

\textbf{Universo.} Corresponde al total de usuarios potenciales del transporte público urbano de la ciudad de Arequipa.

\textbf{Población.} Usuarios con desplazamientos frecuentes en rutas urbanas y acceso a un dispositivo con navegador web y conectividad.

\textbf{Muestra.} Se emplea muestra no probabilística intencional para piloto controlado, conformada por usuarios disponibles en corredores priorizados. Esta decisión es consistente con el carácter aplicado del estudio y con la lógica de validación progresiva del enfoque Scrum.

\subsection*{Marco de trabajo Scrum}

\textbf{A. Roles definidos}
\begin{itemize}
    \item \textbf{Product Owner (investigador):} prioriza el Product Backlog según impacto en reducción de tiempos y valor para el usuario.
    \item \textbf{Scrum Master (investigador/asesor):} asegura cumplimiento del marco Scrum, elimina bloqueos y protege el flujo del Sprint.
    \item \textbf{Equipo de desarrollo:} implementa funcionalidades, integra componentes y ejecuta pruebas técnicas/funcionales.
\end{itemize}

\textbf{B. Artefactos}
\begin{itemize}
    \item \textbf{Product Backlog:} lista priorizada de historias de usuario y requerimientos técnicos.
    \item \textbf{Sprint Backlog:} tareas comprometidas para cada iteración.
    \item \textbf{Incremento:} versión funcional verificable al cierre de cada Sprint.
\end{itemize}

\textbf{C. Eventos}
\begin{itemize}
    \item \textbf{Sprint Planning:} definición de objetivo y alcance del Sprint.
    \item \textbf{Daily Scrum:} seguimiento diario de avances y bloqueos.
    \item \textbf{Sprint Review:} demostración del incremento y validación con criterios de aceptación.
    \item \textbf{Sprint Retrospective:} identificación de mejoras para el siguiente ciclo.
\end{itemize}

\subsubsection*{Plan de trabajo y cronograma (marzo - julio 2026)}

El cronograma se organiza en bloques semanales y quincenales, respetando las fechas establecidas para el proyecto:

\begin{longtable}{|p{0.45\textwidth}|p{0.22\textwidth}|p{0.22\textwidth}|}
\hline
\textbf{Actividad} & \textbf{Fecha de inicio} & \textbf{Fecha de finalización} \\
\hline
\endfirsthead
\hline
\textbf{Actividad} & \textbf{Fecha de inicio} & \textbf{Fecha de finalización} \\
\hline
\endhead
Diseño metodológico & 23/03/2026 & 29/03/2026 \\
\hline
Requerimientos y diseño de la propuesta & 30/03/2026 & 05/04/2026 \\
\hline
Implementación de la propuesta & 06/04/2026 & 12/04/2026 \\
\hline
Validación y recopilación de datos (métricas e instrumentos) & 13/04/2026 & 19/04/2026 \\
\hline
Análisis e interpretación de datos (estadística) & 20/04/2026 & 26/04/2026 \\
\hline
Presentación de resultados y discusión & 27/04/2026 & 03/05/2026 \\
\hline
Formulación de conclusiones y recomendaciones & 04/05/2026 & 10/05/2026 \\
\hline
Aspectos formales (redacción, APA/referencias) & 11/05/2026 & 24/05/2026 \\
\hline
Preparación para la presentación (material gráfico y oratoria) & 25/05/2026 & 05/07/2026 \\
\hline
\end{longtable}

\textbf{Alineación del cronograma en 10 Sprints}
\begin{itemize}
    \item \textbf{Sprint 1 (23/03--29/03).} \textbf{Actividades:} (1) diseño metodológico del estudio; (2) definición de criterios de validación. \textbf{Entregables:} (1) plan metodológico validado; (2) matriz de criterios e indicadores.
    \item \textbf{Sprint 2 (30/03--05/04).} \textbf{Actividades:} (1) levantamiento de requerimientos; (2) priorización del backlog del producto. \textbf{Entregables:} (1) documento de requerimientos; (2) Product Backlog priorizado.
    \item \textbf{Sprint 3 (06/04--12/04).} \textbf{Actividades:} (1) diseño técnico de arquitectura web; (2) diseño de flujo de datos GPS. \textbf{Entregables:} (1) diagrama de arquitectura; (2) especificación de flujo de datos.
    \item \textbf{Sprint 4 (13/04--19/04).} \textbf{Actividades:} (1) desarrollo del módulo de captura/consulta; (2) integración inicial con lógica de tiempos. \textbf{Entregables:} (1) incremento funcional v1; (2) módulo de estimación de llegada operativo.
    \item \textbf{Sprint 5 (20/04--26/04).} \textbf{Actividades:} (1) pruebas funcionales internas; (2) ajustes de estabilidad del sistema. \textbf{Entregables:} (1) reporte de pruebas internas; (2) incremento funcional v2 estabilizado.
    \item \textbf{Sprint 6 (27/04--03/05).} \textbf{Actividades:} (1) validación en entorno piloto; (2) recopilación de métricas pretest--postest. \textbf{Entregables:} (1) base de datos de validación; (2) registro comparativo de tiempos.
    \item \textbf{Sprint 7 (04/05--10/05).} \textbf{Actividades:} (1) análisis estadístico de resultados; (2) elaboración de tablas y gráficos. \textbf{Entregables:} (1) informe estadístico; (2) anexos gráficos de resultados.
    \item \textbf{Sprint 8 (11/05--24/05).} \textbf{Actividades:} (1) redacción de resultados y discusión; (2) formulación de conclusiones y recomendaciones. \textbf{Entregables:} (1) capítulo de resultados y discusión; (2) capítulo de conclusiones y recomendaciones.
    \item \textbf{Sprint 9 (25/05--14/06).} \textbf{Actividades:} (1) normalización formal del documento (APA y referencias); (2) corrección integral de estilo académico. \textbf{Entregables:} (1) tesis final normalizada; (2) versión lista para revisión final.
    \item \textbf{Sprint 10 (15/06--05/07).} \textbf{Actividades:} (1) elaboración de material de sustentación; (2) preparación de oratoria y simulacros. \textbf{Entregables:} (1) presentación final; (2) guion y checklist de defensa.
\end{itemize}

\subsection*{Diseño de investigación y recolección de datos}

Se aplican técnicas mixtas con foco en evidencia empírica del problema:
\begin{itemize}
    \item \textbf{Observación operativa:} registro de tiempos de espera en paraderos y tiempos de desplazamiento.
    \item \textbf{Análisis de registros digitales:} datos GPS y eventos del sistema web.
    \item \textbf{Instrumento pretest--postest:} ficha comparativa de tiempos antes/después de usar la propuesta.
    \item \textbf{Revisión documental técnica:} soporte de decisiones de diseño y validación.
\end{itemize}

Instrumentos principales: prototipo web, base de datos de eventos, fichas de registro temporal y matriz de trazabilidad de incidencias.

\subsection*{Criterios de validación por Sprint}

Al cierre de cada Sprint se verifican criterios de aceptación mínimos:
\begin{itemize}
    \item funcionalidad operativa del incremento;
    \item consistencia de datos capturados;
    \item estabilidad básica de la solución web;
    \item utilidad para la toma de decisión del usuario;
    \item trazabilidad entre requerimiento, implementación y evidencia de validación.
\end{itemize}

\subsection*{Técnicas estadísticas para análisis de datos}

El análisis de resultados se realiza mediante estadística descriptiva y comparación de medidas:
\begin{itemize}
    \item media, mediana y desviación estándar de tiempos de espera y viaje;
    \item variación porcentual pretest--postest para estimar reducción de tiempos;
    \item tablas y gráficos de tendencia por ruta y franja horaria;
    \item interpretación de variaciones para identificar condiciones de mayor impacto.
\end{itemize}

La fase final integra resultados cuantitativos y evaluación funcional para determinar factibilidad de escalamiento de la propuesta.
